%%%%%%%%%%  scan idx 1  |  folio cover-verso  %%%%%%%%%%
\begin{quote}
\textbf{Note de transcription : page encore incomplète.}
Une étiquette collée sur l'exemplaire du scan principal occulte une grande
partie de la notice de l'éditeur. Le témoin fourni commence à la page de titre
et ne contient pas cette page. Les mots occultés restent à établir. Le texte
lisible ci-dessous est conservé; cette note est de la présente transcription.
\end{quote}

\begin{quote}
\textit{Fragments visibles de la partie supérieure de la notice. Les crochets
signalent le texte masqué par l'étiquette, y compris les lettres coupées à son
bord; ils ne précisent pas le nombre de caractères manquants.}
\begin{flushleft}
This series aim[texte masqué]arch and teaching -- quickly, informally\par
and at a high le[texte masqué]ncludes:\par
1. Preliminary d[texte masqué]\par
2. Lectures on a[texte masqué]field\par
3. Seminar work[texte masqué]\par
4. Reports of m[texte masqué]\par
Texts which are[texte masqué]red if they fall within these categories.\par
The timeliness o[texte masqué]which may be unfinished or tentative.\par
Thus, in some i[texte masqué]ts presented which have been or will\par
later be publishe[texte masqué]\par
Manuscripts shou[texte masqué]\par
Publication of \textit{L}[texte masqué]ational mathematical community, in\par
that a commercia[texte masqué]tribution to documents which would\par
otherwise have a[texte masqué]ted, they can be documented in the\par
scientific literatur[texte masqué]
\end{flushleft}
\end{quote}

\subsection*{Manuscripts}
Manuscripts are re[texte masqué]herefore be typed with extreme care.

Symbols not on the typewriter should be inserted by hand in indelible black ink.
Corrections to the typescript should be made by sticking the amended text over
the old one, or by obliterating errors with white correcting fluid. Authors
receive 75 free copies.

The typescript is reduced slightly in size during reproduction; best results
will not be obtained unless the text on any one page is kept within the overall
limit of $18\times26.5$ cm ($7\times10\frac12$ inches). The publishers will be
pleased to supply on request special stationery with the typing area outlined.

Manuscripts in English, German or French should be sent to Prof. Dr. A. Dold,
Mathematisches Institut der Universität Heidelberg, 69 Heidelberg/Germany,
Tiergartenstraße or Prof. Dr. B. Eckmann, Eidgenössische Technische Hochschule,
CH-8006 Zürich/Switzerland.

\subsection*{Lecture Notes in Physics}
\textbf{Bisher erschienen/Already published}

Vol. 1: J. C. Erdmann, Wärmeleitung in Kristallen, theoretische Grundlagen und
fortgeschrittene experimentelle Methoden. II, 283 Seiten. 1969. DM 20,--

Vol. 2: K. Hepp, Théorie de la renormalisation. III, 215 pages. 1969. DM 18,--

Vol. 3: A. Martin, Scattering Theory: Unitarity, Analytic and Crossing.
IV, 125 pages. 1969. DM 14,--

Vol. 4: G. Ludwig, Deutung des Begriffs physikalische Theorie und axiomatische
Grundlegung der Hilbertraumstruktur der Quantenmechanik durch Hauptsätze des
Messens. XI, 469 Seiten. 1970. DM 28,--

Vol. 5: M. Schaaf, The Reduction of the Product of Two Irreducible Unitary
Representations of the Proper Orthochronous Quantummechanical Poincaré Group.
IV, 120 pages. 1970. DM 14,--

Vol. 6: Group Representations in Mathematics and Physics. Edited by V. Bargmann.
V, 340 pages. 1970. DM 24,--

Vol. 7: R. Balescu, J. L. Lebowitz, I. Prigogine, P. Résibois, Z. W. Salsburg,
Lectures in Statistical Physics. V, 181 pages. 1971. DM 18,--

Vol. 8: Proceedings of the Second International Conference on Numerical Methods
in Fluid Dynamics. Edited by M. Holt. IX, 462 pages. 1971. DM 28,--

